% ---------------------------------------------------------------------------
% Cost-aware wall-clock study (Sept 2026 revision). Everything in this file is
% new material and is rendered in red through the enclosing \edit{...}.
% Numbers are pulled from costaware_tables.tex (generated by make_tables.py).
% ---------------------------------------------------------------------------
\section{Cost-Aware Routing and Wall-Clock Evaluation} \label{sec:wallclock}

The experiments of \Cref{sec:experiments} compare methods per iteration, which favors methods whose iterations are expensive. This appendix evaluates the end-to-end \emph{wall-clock time to a target accuracy} of the learned router against HINTS and the classical solver alone, for all five classical solvers and three PDEs, on a $\caN\times\caN$ grid ($N^2 = \caN^2$ unknowns), using the cost-aware instantiation of the greedy rule introduced at the end of \Cref{sec:greedy}. Besides the Poisson and convection--diffusion equations of \Cref{sec:experiments} we add \emph{anisotropic diffusion}, $-\epsilon\,\partial_{x_1}^2 u - \partial_{x_2}^2 u = f$ with $\epsilon = 10^{-2}$ (diffusion tensor $\mathrm{diag}(\epsilon, 1)$), on the same periodic domain with the same forcing distribution. It is the standard model problem on which point-relaxation methods fail to smooth: an error component oscillatory in $x_1$ and smooth in $x_2$ is damped by only $(1-\epsilon)/(1+\epsilon) \approx 0.98$ per Jacobi sweep, and geometric multigrid with point smoothing loses its grid-independent convergence (line relaxation or semi-coarsening would be needed). The same pipeline is then applied to solver ensembles, to larger grids and against stronger classical baselines.

\subsection{Protocol}

\paragraph{Fast, matched implementations.} A wall-clock comparison is meaningful only if the classical baseline is not handicapped by implementation. All solvers are therefore re-implemented for the periodic constant-coefficient setting as $O(N^2)$ operations: Jacobi and damped Jacobi as 5-point stencil sweeps, Gauss--Seidel, SOR and symmetric GS via cached sparse triangular factorizations. Their one-step iterates match the dense implementations used in the rest of the paper to $10^{-12}$ in double precision. Reference solutions are computed by FFT diagonalization of the same discrete operator, i.e., they are exact solutions of the discrete system, so true errors are available for evaluation at no charged cost. All timings are single-threaded CPU measurements (Apple M4 Pro, numpy/scipy/PyTorch, one thread) on an otherwise idle machine.

\paragraph{Corrector.} The neural corrector is a DeepONet applied to the residual $r = f_h - \mathcal{L}_h^a u^{(t)}_h$ and evaluated on a $64\times 64$ sensor grid, i.e., on the $4\times$ coarser grid obtained by band-limited (FFT) restriction to the modes $|k_1|,|k_2| \le 31$ (for the anisotropic equation the sensor grid is $128\times 32$, i.e., full resolution along $x_1$, so that the band contains every mode the point smoothers cannot damp; the sensor grid is the only PDE-specific hyperparameter of the pipeline); its output is prolongated back to the fine grid by exact trigonometric interpolation, so the correction is band-limited by construction and never injects error into the high-frequency band that the classical smoother owns. The trunk net is the fixed orthonormal Fourier basis of that band ($p = 3969$ basis functions, as in POD-DeepONet \citep{lu2022comprehensive}), and the branch net is a linear layer on the $4096$ sensor values; it is fit by least squares on $64{,}000$ exact (residual, error) pairs of the discrete operator drawn from a mixture of spectral shapes (forcing-like, residual-like, white and randomly tilted spectra), and reaches a validation relative error of $1.7\times 10^{-6}$ on every mode of the band. The corrector is applied scale-equivariantly, $C_{\mathrm{NO}}(r) = \|Rr\|\,\mathcal{G}_\theta(Rr/\|Rr\|)$, which preserves $C_{\mathrm{NO}}(0) = 0$ (Proposition~\ref{th:weaklyalphasupermodular}). One corrector call costs $\approx 0.6$--$0.7$\,ms at $128^2$ (two FFTs and a $4096\times 4096$ matrix--vector product), i.e., about $12$ Jacobi sweeps, $3$ GS/SOR sweeps, or $2$ symmetric-GS sweeps (\Cref{tab:ca_costs}). The same corrector is used by HINTS and by all routing policies.

\paragraph{Cost-equalised macro-actions.} Per-iteration costs $c_j$ (residual evaluation plus the update) are measured once per pairing on the benchmark machine (minimum over seven repeated blocks of the block median, cached), the unit of cost $u$ is one corrector call, and every operation is applied $m_j = \max\{1,\mathrm{round}(u/c_j)\}$ times per decision (\Cref{tab:ca_costs}); for the five stationary solvers this equalises costs to within $8\%$. An operation dearer than the unit (a multigrid cycle) is applied once and compared per unit of cost, $\|e_j\|^{u/(m_j c_j)}$; the surrogate weights $\tilde c_j$ use the same exponent. The cost-aware oracle runs \Cref{alg:greedy} on these macro-actions with the true error; the \emph{greedy oracle} row runs the cost-agnostic \Cref{alg:greedy} on single operations, as in the main text.

\paragraph{Lightweight router.} The learned router selects macro-actions from $6 + K$ scalar features, all $O(1)$ given the residual norm that any stopping test computes: the log relative residual, its per-iteration change over the last macro-action and an exponential moving average thereof, the log epoch index, a one-hot encoding of the previous macro-action, and the log counts of corrector calls so far and of iterations since the last call. It is a two-hidden-layer MLP (64 units) evaluated in numpy, costing $\approx 5\,\mu$s per decision independently of $N$. It is trained with the surrogate of \Cref{section:surrogate}, with the per-state weights $\tilde c_k$ normalised by $\sum_k \tilde c_k$ (which does not change the Bayes decision), on states visited by oracle rollouts with $\varepsilon$-exploration (probability $0.1$), followed by two DAgger rounds \citep{ross2011reduction} in which the current router's own rollouts are relabelled by the oracle---the batch analogue of the scheduled sampling used in \Cref{sec:training details}. Training uses $128$ instances per round drawn from the same GRF distribution as the test set but with a different seed; the router of every pairing and ensemble is trained in under two minutes.

\paragraph{Baselines and metric.} HINTS applies the corrector every $\tau$-th iteration; we report $\tau = 25$ (the schedule used in the main text and in \citet{zhang2024blending}) and the best of $\tau \in \{5, 10, 25, 50\}$ selected \emph{per cell} on the test instances themselves (an optimistic baseline). For each of $64$ test instances we record, after every iteration, the true relative error $\|e^{(t)}_h\|/\|u_h\|$ (untimed, benchmark-only) and the cumulative charged wall-clock time of a replay that executes only the chosen operations (residual, update, and---for the learned router---every feature and decision computation; oracle decisions are not charged). We report the median first time at which the error drops below a tolerance $\varepsilon$, and paired per-instance median speedups. We highlight $\varepsilon = h^2 = 6.1\times 10^{-5}$: the discretization error of the second-order scheme scales as $O(h^2)$, so driving the algebraic error far below $h^2$ does not improve the computed solution. We also report the iteration-based metrics of the main text (final relative error and AUC $=\sum_{t=1}^{T}\|e^{(t)}_h\|/\|u_h\|$ over $T = \caT$ iterations).

\subsection{Pairwise results}

\Cref{tab:ca_main} (main text) reports the time to truncation-level accuracy for all ten pairings; the tables below give the comparison against HINTS at three tolerances (\Cref{tab:ca_vshints}), the iteration-based metrics of the main text (\Cref{tab:ca_auc}), corrector-call and iteration counts (\Cref{tab:ca_usage}), the measured costs (\Cref{tab:ca_costs}), and the full tolerance sweeps (\Cref{tab:ca_tol_poisson,tab:ca_tol_convdiff}).

\paragraph{Uniform gains.} In every one of the $\caNumCells$ solver--equation pairings the learned router is the fastest deployable method at $\varepsilon = h^2$: \caSpSolverMin--\caSpSolverMax\ faster than the solver alone, \caSpHintsMin--\caSpHintsMax\ faster than HINTS with $\tau{=}25$, and \caSpBestMin--\caSpBestMax\ faster than the best fixed period chosen per pairing on the test set (all paired Wilcoxon $p < 0.01$). The gains are not confined to the truncation level: at $\varepsilon = 10^{-8}$ the router is still \caSpHintsDeepMin--\caSpHintsDeepMax\ faster than HINTS ($\tau{=}25$) and \caSpBestDeepMin--\caSpBestDeepMax\ faster than the per-tolerance best period (\Cref{tab:ca_vshints}), and the best period itself varies across pairings and tolerances ($\tau{=}5$ at truncation level, $\tau \in \{5, 10, 50\}$ at $10^{-8}$), so no single fixed schedule is competitive everywhere. The learned router reproduces the cost-aware oracle's decisions almost exactly (the same median iteration and call counts in \Cref{tab:ca_usage}; median times within \caSpOracleRatioMin--\caSpOracleRatioMax\ of the oracle, the difference being its own decision cost and timer noise).

\paragraph{Where the time goes.} \Cref{tab:ca_usage} explains the numbers: with a corrector that removes the smooth band to $10^{-6}$ in one call, the cost-aware oracle and the router call it first and reach $h^2$ after one or two further sweeps, i.e., in about two iterations, whereas HINTS reaches the same accuracy only at its first scheduled call (iteration $\tau$), having spent $\tau{-}1$ sweeps on an error that is almost entirely smooth. Below $h^2$ the remaining error is the high-frequency band that the corrector cannot see; the router therefore stops calling it (one to three calls in total down to $10^{-8}$; \Cref{fig:ca_usage}), whereas a fixed schedule keeps paying for a call every $\tau$ iterations---HINTS ($\tau{=}5$) executes hundreds of useless calls on the Jacobi pairings (\Cref{tab:ca_tol_poisson}). The cost-agnostic greedy oracle (\Cref{alg:greedy} on single operations) makes the same decisions as its cost-aware counterpart here, because after the first call the corrector is never better than a single sweep on the remaining high-frequency error; the two rules differ only through the granularity of their decisions, and the learned router pays for neither.

\paragraph{Anisotropic diffusion.} The anisotropic equation is where the two components of the hybrid are most complementary and where the classical baselines are weakest: the point smoothers cannot damp the $x_1$-oscillatory modes at all (\Cref{tab:ca_tol_aniso}: none of the five solvers alone reaches $h^2$ within $60{,}000$ iterations), whereas the corrector, whose sensor grid keeps full resolution along $x_1$, removes exactly those modes together with the smooth ones in a single call and leaves the smoother only the $x_2$-oscillatory band, which it damps at the isotropic rate. Consequently the learned router reaches $h^2$ in about $1$\,ms for every pairing---the same absolute time as on the isotropic equations---while HINTS ($\tau{=}25$) needs $3$--$10$\,ms and the best fixed period $1.2$--$2.4$\,ms (\Cref{tab:ca_main,tab:ca_vshints}); the router again matches the cost-aware oracle. The multigrid and Krylov baselines are not run for this equation: with point smoothing, the V-cycle's contraction factor degrades from $\approx 0.03$ to $\approx 0.85$ per cycle (line relaxation or semi-coarsening would be required), so the comparison of \Cref{sec:wallclock_strong} would not be informative.

\paragraph{Iteration-based metrics.} \Cref{tab:ca_auc} confirms that the iteration-count conclusions of \Cref{sec:experiments} hold on the finer grid: the router's AUC over $T = \caT$ iterations is three to four orders of magnitude below HINTS' and four to five below the solver-only baseline (all $p < 10^{-3}$), and its error after $T$ iterations is at the stopping threshold of the runs ($10^{-9}$) except for undamped Jacobi, where all hybrids stall at the near-Nyquist band.


\begin{table}[H]
\color{red}
\centering
\caption{Paired median speedup of the learned router over HINTS at three tolerances ($64$ instances; bold: $\ge 1.1\times$ with one-sided Wilcoxon $p < 0.01$). The right block compares against the fixed period that is best \emph{for that tolerance and pairing}, chosen on the test set.}
\resizebox{0.95\textwidth}{!}{\cavshints}
\label{tab:ca_vshints}
\end{table}

\begin{table}[H]
\color{red}
\centering
\caption{Iteration-based metrics of the main text on the $\caN\times\caN$ grid: relative error after $T = \caT$ iterations and AUC of the relative error over $T$ iterations, mean (standard error) over $64$ instances; $p$-values are one-sided paired $t$-tests of the baseline's AUC against the learned router's.}
\resizebox{\textwidth}{!}{\caauc}
\label{tab:ca_auc}
\end{table}

\begin{table}[H]
\color{red}
\centering
\caption{Median number of iterations and corrector calls needed to reach $\varepsilon = h^2$ on the $\caN\times\caN$ grid.}
\resizebox{0.8\textwidth}{!}{\causage}
\label{tab:ca_usage}
\end{table}

\begin{table}[H]
\color{red}
\centering
\caption{Measured single-thread per-iteration costs (residual + update) on the $\caN\times\caN$ grid and the resulting macro-action sizes $m_j = \max\{1,\mathrm{round}(u/c_j)\}$ ($u$: corrector call).}
\resizebox{0.75\textwidth}{!}{\cacosts}
\label{tab:ca_costs}
\end{table}

\begin{table}[H]
\color{red}
\centering
\caption{Poisson on the $\caN\times\caN$ grid: median wall-clock time to relative error $\varepsilon$ for every policy and pairing ($64$ instances; paired median speedup over solver-only in parentheses; $\dagger n$: $n$ instances did not reach the tolerance within the iteration cap).}
\resizebox{\textwidth}{!}{\cawctolpoisson}
\label{tab:ca_tol_poisson}
\end{table}

\begin{table}[H]
\color{red}
\centering
\caption{Convection--diffusion on the $\caN\times\caN$ grid: median wall-clock time to relative error $\varepsilon$ for every policy and pairing ($64$ instances).}
\resizebox{\textwidth}{!}{\cawctolconv}
\label{tab:ca_tol_convdiff}
\end{table}

\begin{table}[H]
\color{red}
\centering
\caption{Anisotropic diffusion on the $\caN\times\caN$ grid: median wall-clock time to relative error $\varepsilon$ for every policy and pairing ($64$ instances). The point smoothers alone do not reach $h^2$ within $60{,}000$ iterations (their times are lower bounds).}
\resizebox{\textwidth}{!}{\cawctolaniso}
\label{tab:ca_tol_aniso}
\end{table}

\begin{figure}[H]
\color{red}
    \centering
    \newfig[0.93\linewidth]{neurips_images/ca_deeponet_predictions.png}
    \caption{One-shot predictions of the corrector on two test forcings per equation (rows): forcing, exact discrete solution, DeepONet output $C_{\mathrm{NO}}(f)$ and its error. The corrector solves the band $|k|\le 31$ to $\approx 10^{-6}$; the remaining error is the high-frequency content of the solution outside the band, which the classical smoother removes in a few sweeps.}
    \label{fig:ca_predictions}
\end{figure}

\begin{figure}[H]
\color{red}
    \centering
    \newfig[0.98\linewidth]{neurips_images/ca_router_usage.png}
    \caption{Corrector usage: fraction of test runs that execute a corrector call at each iteration, for the learned router, the cost-aware oracle and HINTS ($\tau = 25$). The router calls the corrector once at the start, spends a solver-dependent number of sweeps on the high-frequency band, and calls it again only when the smooth error has re-emerged; HINTS waits $24$ sweeps before its first call.}
    \label{fig:ca_usage}
\end{figure}

\begin{figure}[H]
\color{red}
    \centering
    \newfig[0.98\linewidth]{neurips_images/ca_router_decisions_poisson.png}
    \newfig[0.98\linewidth]{neurips_images/ca_router_decisions_convdiff.png}
    \caption{Routing decisions (dark: corrector call) of the learned router (top rows) and the cost-aware oracle (bottom rows) over the first $60$ iterations of the stored test instances, for every pairing.}
    \label{fig:ca_decisions}
\end{figure}

\begin{figure}[H]
\color{red}
    \centering
    \newfig[0.98\linewidth]{neurips_images/ca_convergence.png}
    \caption{Relative error against charged wall-clock time for one test instance per pairing. The dashed horizontal line is $\varepsilon = h^2$.}
    \label{fig:ca_convergence}
\end{figure}


\subsection{Solver ensembles}

The ensemble experiments of \Cref{sec:experiments} are repeated in wall-clock terms with the cost-aware router over $\mathrm{NO}\cup\mathcal{W}$ for the same four ensembles $\mathcal{W}$ (macro-action sizes are set from the measured costs of all members, $m_j = \mathrm{round}(c_{\max}/c_j)$). HINTS does not apply to ensembles, so the comparison is against every member solver alone and against the best pairwise router $\mathrm{Router}(\mathrm{NO}\cup\{\text{solver}\})$, $\text{solver}\in\mathcal{W}$, from \Cref{tab:ca_main}.

\paragraph{Results.} \Cref{tab:ca_ens} shows that the ensemble router is competitive with, but does not systematically improve on, the best pairwise router of its members: across the $\caNumEns$ ensemble--equation cells the ratio of median times to $h^2$ (best pairwise router over ensemble router) ranges from \caEnsVsPairMin\ to \caEnsVsPairMax, i.e., within $\pm 15\%$, while the ensemble router is \caEnsVsSolverMin--\caEnsVsSolverMax\ faster than the fastest member solver alone, and it does not need to know in advance which member is fastest (the best pairwise member differs between Jacobi and damped Jacobi across cells). The cost-aware \emph{oracle} over the full ensemble is faster than every pairwise router in every cell (\Cref{tab:ca_ens}), so the ensembles do contain exploitable diversity; the gap between the ensemble router and its oracle (significant in the larger ensembles, \Cref{tab:ca_stats_ens}) is imitation error. It has a clear cause: with several cost-equalised smoothers in the ensemble the oracle's choice among them depends on the spectral content of the current error, which the router's residual-norm history features only reflect indirectly, and the surrogate's weights are small for such near-ties, so the router agrees with the oracle on only $55$--$83\%$ of its own decisions (against $90$--$99\%$ for the pairwise routers) and occasionally picks a smoother that is a few percent worse per unit cost. Deciding \emph{when} to call the corrector---the decision that dominates the wall-clock gains---is learned equally well in the ensembles. \Cref{tab:ca_ensusage} lists how the iterations up to the $h^2$ crossing are distributed: the corrector is called once (or twice) and the remaining iterations go to whichever smoother the router selects for the high-frequency band, with an instance-dependent mix, as in \Cref{sec:solverdecisions}.

\begin{table}[H]
\color{red}
\centering
\caption{Ensembles: median wall-clock time to $\varepsilon = h^2$ on the $\caN\times\caN$ grid ($64$ instances). The ensemble router is compared with the fastest member solver alone, with the fastest pairwise router of its members (paired median speedup in parentheses), and with the cost-aware oracle over the same ensemble.}
\resizebox{\textwidth}{!}{\caens}
\label{tab:ca_ens}
\end{table}

\begin{table}[H]
\color{red}
\centering
\caption{Ensembles: average fraction of iterations (up to the $h^2$ crossing) spent in each component by the learned router, mean (standard deviation) over $64$ instances; ``-'' marks components not in the ensemble.}
\resizebox{\textwidth}{!}{\caensusage}
\label{tab:ca_ensusage}
\end{table}

\subsection{Stronger classical baselines, larger grids and overheads}
\label{sec:wallclock_strong}

The comparisons above are against the stationary solvers of \Cref{sec:experiments}. This section adds the classical methods a practitioner would actually reach for on these problems, scales the grid up to $512\times 512$, and accounts explicitly for every overhead.

\paragraph{Baselines.} All baselines run on the same $O(N^2)$ stencil with the same accounting (first iteration at which the true error is below $\varepsilon$; timed replay of exactly that many iterations). (i) \emph{FFT direct solve}: for the constant-coefficient periodic problems used here the discrete operator is diagonalised by the FFT, so the system can be solved directly with two FFTs; this is the reference solver we use to compute ground truth and it is included as the lower bound on the time any iterative method can achieve on this problem class. (ii) \emph{Geometric multigrid}: V(2,2)-cycles with lexicographic Gauss--Seidel smoothing, full-weighting restriction, bilinear prolongation, rediscretised coarse operators and a direct solve on the $32\times 32$ coarsest grid (error contraction $\approx 0.03$ per cycle). A V-cycle is a stationary iteration, so multigrid is also used as an ensemble member below. (iii) \emph{Krylov methods} (SciPy implementations driven by the stencil matvec): conjugate gradients, CG preconditioned by symmetric GS and by a symmetric V-cycle for Poisson; BiCGSTAB, multigrid-preconditioned BiCGSTAB and restarted GMRES(20) for convection--diffusion (GMRES is accounted at the granularity of its restart cycles).

\begin{table}[H]
\color{red}
\centering
\caption{Strong classical baselines on the $128\times 128$ grid: median time to relative error $\varepsilon$ over $64$ instances (paired median speedup of the learned router with its best stationary pairing over each method in parentheses; $>1$ means the router is faster) and median iterations to $h^2$. The FFT solve is a direct method for this problem class and bounds every iterative solver from below.}
\resizebox{\textwidth}{!}{\cabaselines}
\label{tab:ca_baselines}
\end{table}

\paragraph{Results at $128^2$.} With the same implementation technology (numpy stencils and scipy sparse triangular solves for every classical operation; a BLAS matrix--vector product for the corrector), the learned router with its best stationary pairing reaches truncation-level accuracy \caVsMgMin--\caVsMgMax\ faster than multigrid alone (three V-cycles), \caVsKrylovMin--\caVsKrylovMax\ faster than the best multigrid-preconditioned Krylov method and an order of magnitude faster than the unpreconditioned or SymGS-preconditioned Krylov solvers; the only faster method is the FFT direct solve, by \caFftRatioMin--\caFftRatioMax, which exists only because the test problems have constant coefficients and periodic boundaries (it is not an iterative method and is the reference we solve against). The reason the hybrid beats multigrid here is the corrector's accuracy on the smooth band: one call removes the modes $|k|\le 31$ to $10^{-6}$ at the price of a single $4096\times4096$ matrix--vector product, whereas a V(2,2)-cycle contracts the whole spectrum by $\approx 0.03$ at the price of four Gauss--Seidel sweeps on the fine grid plus the coarse-grid work. \Cref{tab:ca_baselines} also reports the router over $\{\mathrm{NO}, \text{multigrid}\}$: the cost-aware oracle and the learned router call the corrector once and then let the V-cycle remove the rest, and are \caVsMgEnsMin--\caVsMgEnsMax\ faster than multigrid alone, i.e., adding a strong classical member to the ensemble is handled by the same construction without any change (the multigrid cycle is dearer than the unit, so it is the case in which the per-unit-cost comparison of \Cref{sec:greedy} matters).

\paragraph{Scaling with the grid.} \Cref{tab:ca_scaling} repeats the Jacobi and Gauss--Seidel pairings on $256\times 256$ and $512\times 512$ grids ($32$ and $16$ instances; correctors retrained per grid with the same recipe, keeping the sensor grid at $64\times 64$, i.e., the corrector always handles the modes $|k|\le 31$, and the router retrained per pairing). The corrector's cost is dominated by a $4096\times 4096$ matrix--vector product and is therefore nearly grid-independent, while every classical iteration grows as $N^2$, so the corrector becomes relatively cheaper as $N$ grows; the classical member, however, has to cover a growing fraction of the spectrum. The solver-only baselines at $512^2$ are capped at $8000$ iterations and reported as lower bounds; restarted GMRES, which was two orders of magnitude slower than the router at $128^2$ and $256^2$, is not run at $512^2$.

\paragraph{What the scaling shows.} At $256^2$ the router remains the fastest iterative method for every pairing and both equations: $1.8$--$10.9\times$ faster than HINTS ($\tau{=}25$), $1.25$--$2.4\times$ faster than the per-cell best period, $4.6$--$6.7\times$ faster than multigrid alone and $5$--$6\times$ faster than multigrid-preconditioned CG/BiCGSTAB (\Cref{tab:ca_main_256,tab:ca_baselines_256}); the FFT direct solve is $1.7$--$2.6\times$ faster than the router. At $512^2$ the picture splits by pairing (\Cref{tab:ca_main_512}). With SOR and multigrid the router is still $3.1\times$ and $2.1\times$ faster than the best fixed schedule and $1.3$--$2.4\times$ faster than multigrid alone; the $\{\mathrm{NO}, \text{multigrid}\}$ router (39\,ms on Poisson) is the fastest iterative method we measured at this size. With Jacobi and Gauss--Seidel, however, the corrector's fixed band $|k|\le 31$ leaves the smoother all modes $k \ge 32$, whose slowest member contracts by only $0.98$ per Jacobi sweep at $N=512$ (against $0.71$ at $N=128$): every policy then needs the same $\approx 80$ sweeps after a single corrector call, the wall-clock time is dominated by sweeps, and the schedule barely matters---the router is $1.11\times$ faster than HINTS on Poisson (GS and Jacobi) but only on par with a fixed $\tau{=}50$ schedule on convection--diffusion with Jacobi, where its imitation of the oracle also degrades (4 of 16 instances do not reach $h^2$ within the cap; the oracle reaches all 16 in the same $84$ iterations as HINTS). This is the expected limit of a corrector whose band does not grow with the grid: the hybrid's advantage comes from removing the modes the smoother cannot damp, and at $512^2$ a $64\times 64$ sensor grid no longer covers them. A corrector with a band that scales with $N$ (at the price of a dense matrix--vector product that grows like the fourth power of the band) or a multigrid member in the ensemble restores the gains, as the multigrid pairing shows.

\begin{table}[H]
\color{red}
\centering
\caption{Scaling: median time to $\varepsilon = h^2$ for the Jacobi and GS pairings against the strongest classical alternatives on each grid (paired median speedup of the learned router over each method in parentheses; $\dagger n$: $n$ instances did not reach the tolerance within the iteration cap).}
\resizebox{\textwidth}{!}{\cascaling}
\label{tab:ca_scaling}
\end{table}

\paragraph{Overheads and amortisation.} \Cref{tab:ca_overheads} lists the measured cost of every operation as a function of $N$, including one decision of the LSTM router of \Cref{sec:experiments} (full-state input of size $4N^2$, hidden size 256, 4 layers): at $128^2$ one LSTM decision costs \caLstmMs\,ms, about \caLstmOverJacobi\ Jacobi sweeps, which is why the wall-clock study uses the scalar-feature router ($\approx 5\,\mu$s, independent of $N$). The learned router's decisions therefore account for well under $1\%$ of its solve time, and its corrector calls for $50$--$70\%$ of the time to $h^2$ (one or two calls). \Cref{tab:ca_amort} accounts for training: the corrector (shared with HINTS) is fit in minutes, the router in seconds to a few minutes per pairing, and the break-even number of solves after which the router's own training cost is recovered relative to HINTS ($\tau{=}25$) with the same corrector.

\begin{table}[H]
\color{red}
\centering
\caption{Measured single-thread cost of one operation as a function of the grid size (Poisson; median of repeated calls). Multigrid uses a direct solve below $64^2$ and is not listed there.}
\resizebox{\textwidth}{!}{\caoverheads}
\label{tab:ca_overheads}
\end{table}

\begin{table}[H]
\color{red}
\centering
\caption{Training-cost amortisation for the GS pairing: corrector training (data generation and least-squares fit, CPU), router training (oracle rollouts and two DAgger rounds), median wall-clock saved per solve at $\varepsilon = h^2$ relative to HINTS ($\tau{=}25$) and to the solver alone, and the number of solves after which the router's training cost is recovered relative to HINTS.}
\resizebox{\textwidth}{!}{\caamort}
\label{tab:ca_amort}
\end{table}

\subsection{Complete results on the $256\times256$ and $512\times512$ grids}
\label{sec:wallclock_large}

The tables of \Cref{sec:wallclock_strong} summarise the larger grids at $\varepsilon = h^2$ for two pairings; the following tables repeat the full per-tolerance results of the $128^2$ study (\Cref{tab:ca_main,tab:ca_vshints,tab:ca_usage,tab:ca_costs,tab:ca_tol_poisson,tab:ca_tol_convdiff}) on the $256\times 256$ grid ($32$ test instances, all six pairings) and on the $512\times 512$ grid ($16$ instances; Jacobi, GS, SOR and multigrid pairings; solver-only runs capped at $8000$ iterations and reported as lower bounds). Correctors and routers are retrained per grid with the recipe of \Cref{sec:wallclock}; the corrector's sensor grid stays at $64\times 64$ (band $|k|\le 31$).

\begin{table}[H]
\color{red}
\centering
\caption{$256\times256$: median wall-clock time to $\varepsilon = h^2 = 1.5\times10^{-5}$ ($32$ instances; paired median speedup over solver-only in parentheses; $\dagger n$: $n$ instances did not reach the tolerance within the $40{,}000$-iteration cap, in which case the median time at the cap is given as a lower bound).}
\resizebox{\textwidth}{!}{\cawcmainB}
\label{tab:ca_main_256}
\end{table}

\begin{table}[H]
\color{red}
\centering
\caption{$256\times256$: paired median speedup of the learned router over HINTS ($32$ instances; bold: $\ge1.1\times$ with one-sided Wilcoxon $p<0.01$).}
\resizebox{0.95\textwidth}{!}{\cavshintsB}
\label{tab:ca_vshints_256}
\end{table}

\begin{table}[H]
\color{red}
\centering
\caption{$256\times256$: median number of iterations and corrector calls to reach $\varepsilon = h^2$.}
\resizebox{0.8\textwidth}{!}{\causageB}
\label{tab:ca_usage_256}
\end{table}

\begin{table}[H]
\color{red}
\centering
\caption{$256\times256$: measured single-thread per-iteration costs and macro-action sizes.}
\resizebox{0.75\textwidth}{!}{\cacostsB}
\label{tab:ca_costs_256}
\end{table}

\begin{table}[H]
\color{red}
\centering
\caption{$256\times256$, Poisson: median wall-clock time to relative error $\varepsilon$ for every policy and pairing ($32$ instances).}
\resizebox{\textwidth}{!}{\cawctolpoissonB}
\label{tab:ca_tol_poisson_256}
\end{table}

\begin{table}[H]
\color{red}
\centering
\caption{$256\times256$, convection--diffusion: median wall-clock time to relative error $\varepsilon$ for every policy and pairing ($32$ instances).}
\resizebox{\textwidth}{!}{\cawctolconvB}
\label{tab:ca_tol_convdiff_256}
\end{table}

\begin{table}[H]
\color{red}
\centering
\caption{$256\times256$, anisotropic diffusion: median wall-clock time to relative error $\varepsilon$ for the Jacobi, GS and SOR pairings ($16$ instances).}
\resizebox{\textwidth}{!}{\cawctolanisoB}
\label{tab:ca_tol_aniso_256}
\end{table}

\begin{table}[H]
\color{red}
\centering
\caption{$256\times256$: strong classical baselines ($32$ instances; paired median speedup of the learned router with its best stationary pairing over each method in parentheses).}
\resizebox{\textwidth}{!}{\cabaselinesB}
\label{tab:ca_baselines_256}
\end{table}

\begin{table}[H]
\color{red}
\centering
\caption{$512\times512$: median wall-clock time to $\varepsilon = h^2 = 3.8\times10^{-6}$ ($16$ instances; Jacobi, GS, SOR and multigrid pairings; solver-only runs capped at $8000$ iterations).}
\resizebox{\textwidth}{!}{\cawcmainC}
\label{tab:ca_main_512}
\end{table}

\begin{table}[H]
\color{red}
\centering
\caption{$512\times512$: paired median speedup of the learned router over HINTS ($16$ instances; bold: $\ge1.1\times$ with $p<0.01$).}
\resizebox{0.95\textwidth}{!}{\cavshintsC}
\label{tab:ca_vshints_512}
\end{table}

\begin{table}[H]
\color{red}
\centering
\caption{$512\times512$: median iterations and corrector calls to $h^2$ (left) and measured per-iteration costs (right).}
\resizebox{0.48\textwidth}{!}{\causageC}\hfill\resizebox{0.48\textwidth}{!}{\cacostsC}
\label{tab:ca_usage_512}
\end{table}

\begin{table}[H]
\color{red}
\centering
\caption{$512\times512$, Poisson: median wall-clock time to relative error $\varepsilon$ for every policy and pairing ($16$ instances).}
\resizebox{\textwidth}{!}{\cawctolpoissonC}
\label{tab:ca_tol_poisson_512}
\end{table}

\begin{table}[H]
\color{red}
\centering
\caption{$512\times512$, convection--diffusion: median wall-clock time to relative error $\varepsilon$ for every policy and pairing ($16$ instances).}
\resizebox{\textwidth}{!}{\cawctolconvC}
\label{tab:ca_tol_convdiff_512}
\end{table}

\begin{table}[H]
\color{red}
\centering
\caption{$512\times512$: strong classical baselines ($16$ instances).}
\resizebox{\textwidth}{!}{\cabaselinesC}
\label{tab:ca_baselines_512}
\end{table}

\subsection{Statistical significance and multi-trial robustness}
\label{sec:wallclock_stats}

All wall-clock comparisons are paired: every policy solves the same $64$ test instances, so we test the per-instance log time ratios with a one-sided Wilcoxon signed-rank test (alternative: the router is faster) and, as a parametric complement, a one-sided paired $t$-test on log times; instances that a baseline fails to bring below the tolerance within the iteration cap are counted against the router (ratio $0$). \Cref{tab:ca_stats} reports the $p$-values of the pairwise comparisons of \Cref{tab:ca_main,tab:ca_vshints}, together with a multi-trial analysis: the router of every pairing was retrained from scratch with five independent seeds (different training instances, exploration and initialisation; the corrector and the test set are fixed) and re-run on the test set. Because these trials ran in separate sessions, whose live timings are not comparable with the main run's, they are compared in \emph{work units} (measured per-iteration costs times executed operations, plus $5\,\mu$s per decision; the main run's HINTS work units serve as the reference), which are determined by the decision sequence alone. We report the mean and standard deviation over seeds of the median work-unit time to $h^2$, in parentheses the fraction of test instances on which the seed router's iteration count to $h^2$ coincides with the main router's (identical decision sequences up to the crossing), the range over seeds of the paired speedup over HINTS ($\tau{=}25$), and the number of seeds for which the router is significantly faster ($p<0.01$) than both HINTS ($\tau{=}25$) and the best fixed period. The learned policy turns out to be insensitive to the training seed: the seed routers reproduce the main router's decision sequence up to the $h^2$ crossing on essentially every test instance, so their work-unit times and speedups coincide and every seed passes both significance tests. \Cref{tab:ca_stats_ens} gives the corresponding tests for the ensembles (including the tests in the opposite direction, to show where an ensemble router is significantly \emph{slower} than the best pairwise router or than its oracle), and \Cref{tab:ca_stats_base} for the strong classical baselines of \Cref{sec:wallclock_strong}. The iteration-based table (\Cref{tab:ca_auc}) already carries paired $t$-tests on the AUC.

\begin{table}[H]
\color{red}
\centering
\caption{Significance of the pairwise comparisons on the $128\times 128$ grid ($64$ paired instances; one-sided Wilcoxon / paired $t$-test $p$-values on live times, ``$10^{-k}$'' denotes $p<10^{-k}$) and robustness over five independent router-training seeds (work-unit times; the percentage is the share of instances whose decision sequence to $h^2$ is identical to the main router's).}
\resizebox{\textwidth}{!}{\castats}
\label{tab:ca_stats}
\end{table}

\begin{table}[H]
\color{red}
\centering
\caption{Ensembles: one-sided Wilcoxon $p$-values at $\varepsilon = h^2$ ($64$ paired instances) for the ensemble router against the fastest member solver alone, against the best pairwise router in both directions, and against the cost-aware oracle over the same ensemble (alternative: the ensemble router is slower).}
\resizebox{\textwidth}{!}{\castatsens}
\label{tab:ca_stats_ens}
\end{table}

\begin{table}[H]
\color{red}
\centering
\caption{Strong baselines: one-sided Wilcoxon $p$-values that the learned router (best stationary pairing) is faster than each method at each tolerance; where the baseline's median is lower the entry gives the $p$-value that the router is slower.}
\resizebox{0.8\textwidth}{!}{\castatsbase}
\label{tab:ca_stats_base}
\end{table}

